\section{Discussion}

MelanomaNet reaches competitive accuracy while exposing four distinct views of
each prediction, and the two goals do not appear to trade off against one
another: the interpretability components are post-hoc and leave the classifier
unchanged. What the components add is a way to check a prediction rather than
accept it. The ABCDE pipeline restates the image in the vocabulary a
dermatologist already uses, so a reviewer can compare the model's stated
features against their own reading; the alignment and faithfulness results give
some assurance that the attention maps track the model rather than decorate it;
and the concept sensitivities describe a decision in semantic terms that
complement the spatial attention map.

The uncertainty estimates are useful precisely when the classifier is wrong but
confident. Separating the epistemic and aleatoric contributions distinguishes
two situations a clinician would treat differently: a case unlike anything in
the training data, where more data would help, and a case that is genuinely
ambiguous in the image itself, where it would not. Flagging such cases for
review is the mechanism by which the system is meant to support a clinician
rather than replace one.

Several limitations bound these claims. The ABCDE analysis depends on automated
segmentation, which degrades on low-contrast lesions and cluttered backgrounds
and propagates its errors into the derived scores. The evolution criterion is
absent because it requires longitudinal images that single-image datasets do
not contain. Class imbalance still limits the rarer classes---actinic keratosis
in particular---despite the weighted loss. The evaluation uses a single
institution-agnostic dataset with no external test set, so the numbers should be
read as internal validation rather than evidence of generalisation. Finally,
post-hoc explanations describe correlations in the trained model; they do not
prove that its internal computation follows clinical reasoning, and should be
presented to clinicians with that caveat.

Future work follows from these limitations: an external validation set from a
different source to test generalisation, longitudinal data to add the evolution
criterion, a larger concept vocabulary, and a reader study measuring whether the
explanations actually change clinician decisions.
